% 02_problem.tex — Rail Paper: The Problem
% ============================================================

\section{The Problem: Solver-Only Optimization}
\label{sec:problem}

The computational imaging community has invested enormous effort in
building better reconstruction algorithms---deeper networks, more
sophisticated architectures, larger training sets---while treating the
forward measurement operator as a fixed, known quantity.
This section demonstrates that this assumption is catastrophically
wrong, and that the resulting \emph{solver-only optimization} paradigm
produces an illusion of progress that collapses on contact with
physical reality.

% ------------------------------------------------------------------
%  Mask-sensitivity spectrum
% ------------------------------------------------------------------

\subsection{The Mask-Sensitivity Spectrum}
\label{sec:sensitivity}

To quantify the scale of the problem, we evaluate five representative
CASSI reconstruction methods under a standardized 4-scenario protocol.
In \textbf{Scenario~I (Ideal)}, the true forward operator $\mathbf{H}$
is used for both measurement simulation and reconstruction---the oracle
upper bound reported in virtually all published work.
In \textbf{Scenario~II (Assumed/Mismatch)}, measurements are generated
with $\mathbf{H}$ but reconstruction uses the nominal operator
$\mathbf{H}_{\mathrm{nom}}$, which differs by a 1-pixel mask shift and
$0.1^\circ$ rotation---a perturbation well within typical alignment
tolerances.
In \textbf{Scenario~III (Corrected)}, a calibrated operator
$\hat{\mathbf{H}}$ is used for reconstruction, obtained via a
lightweight calibration procedure.
The \textbf{Degradation} column reports the drop from Scenario~I to
Scenario~II; the \textbf{Calibration Gain} column reports the
improvement from Scenario~II to Scenario~III.

\begin{table}[t]
\centering
\caption{%
  \textbf{Mask-sensitivity spectrum for CASSI reconstruction methods.}
  Ideal PSNR (Scenario~I), Mismatch PSNR (Scenario~II), Degradation
  (I$\to$II), and Calibration Gain (II$\to$III).
  All values in dB, averaged over 10 test scenes.
  Neural solvers (MST-S, MST-L, HDNet) exhibit dramatically larger
  mismatch degradation than classical methods (GAP-TV, PnP-HSICNN),
  confirming that model capacity amplifies operator sensitivity.%
}
\label{tab:sensitivity}
\smallskip
\begin{tabular}{@{}lcccc@{}}
\toprule
\textbf{Method}
  & \textbf{Ideal (I)}
  & \textbf{Mismatch (II)}
  & \textbf{Degradation}
  & \textbf{Calib.\ Gain} \\
\midrule
GAP-TV~\cite{yuan2016gaptv}
  & 24.22  & 19.66  & \worst{$-4.56$}   & $+0.51$ \\
PnP-HSICNN~\cite{wang2020pnp}
  & 25.12  & 19.10  & \worst{$-6.02$}   & $+0.71$ \\
MST-S~\cite{cai2022mst}
  & 33.98  & 18.01  & \worst{$-15.97$}  & \best{$+3.00$} \\
MST-L~\cite{cai2022mst}
  & \best{34.81}  & 18.09  & \worst{$-16.72$}  & \best{$+3.01$} \\
HDNet~\cite{hu2022hdnet}
  & 34.66  & \best{24.18}  & \worst{$-10.47$}  & $+0.05$ \\
\bottomrule
\end{tabular}
\end{table}

The results in \cref{tab:sensitivity} reveal a striking pattern.
Under ideal conditions, MST-L leads the field at 34.81\,dB---a full
10.59\,dB above GAP-TV.
Under mismatch, MST-L collapses to 18.09\,dB, while GAP-TV drops only
to 19.66\,dB.
The classical method \emph{outperforms} the state-of-the-art
transformer under realistic physical conditions.
\Cref{fig:sensitivity} visualizes this inversion as a bar chart,
making the crossover unmistakable.

\paragraph{Key insight: neural solvers amplify mismatch sensitivity.}
The degradation column tells the story most clearly:
GAP-TV loses 4.56\,dB, PnP-HSICNN loses 6.02\,dB, and MST-L loses
16.72\,dB.
There is a near-perfect inverse relationship between ideal-condition
performance and mismatch robustness.
The same representational capacity that enables MST-L to learn subtle
spectral correlations in the training data also enables it to overfit
the precise structure of the forward operator it was trained with.
When that operator changes by even a single pixel, the learned features
become not merely useless but actively harmful, producing structured
artifacts that are worse than the blurry but stable output of GAP-TV.

HDNet presents a partial counterexample: it achieves high ideal PSNR
(34.66\,dB) with smaller degradation ($-10.47$\,dB), retaining
24.18\,dB under mismatch.
However, its calibration gain is negligible ($+0.05$\,dB), suggesting
that its mask-conditioning pathway has already saturated what it can
extract from the nominal operator.
This taxonomy of solver behavior---mask-oblivious, mask-conditioned, and
mask-adapted---has important implications for the design of
calibration-aware architectures.

% ------------------------------------------------------------------
%  Evaluation fragmentation
% ------------------------------------------------------------------

\subsection{Evaluation Fragmentation}
\label{sec:fragmentation}

The mask-sensitivity spectrum is only one dimension of a much larger
problem.
Even if every lab evaluated mismatch sensitivity, the results would
remain incomparable due to pervasive \emph{evaluation fragmentation}:

\begin{itemize}
  \item \textbf{Metrics.}
        Some groups report peak signal-to-noise ratio (PSNR); others
        report structural similarity (SSIM), spectral angle mapper
        (SAM), or learned perceptual image patch similarity (LPIPS).
        Metrics are computed on different dynamic ranges, with or
        without border cropping, and averaged over different test splits.

  \item \textbf{Datasets.}
        CASSI methods are evaluated on KAIST~\cite{choi2017kaist},
        CAVE~\cite{yasuma2010cave}, ICVL~\cite{arad2016icvl}, or
        custom datasets with different spatial resolutions, spectral
        bands, and scene statistics.
        Results on one dataset do not transfer to another.

  \item \textbf{Noise models.}
        Some methods assume noiseless measurements; others add Gaussian
        noise at varying SNRs; still others use Poisson-Gaussian models.
        The noise model interacts with mismatch in complex ways that
        make cross-paper comparisons meaningless.

  \item \textbf{Mismatch models.}
        The few papers that consider operator mismatch define it
        differently: some shift the mask, some perturb dispersion
        parameters, some add calibration error.
        There is no standard mismatch taxonomy or severity scale.
\end{itemize}

The result is a literature in which every paper reports numbers that
cannot be meaningfully compared with any other paper.
Progress appears rapid when measured within a single group's evaluation
setup, but the aggregate state of the field is unknown.
This is precisely the condition that the SolveEverything framework
calls \emph{The Muddle}~\cite{solveeverything}: everyone is working
hard, but no one can tell whether the work is adding up.

% ------------------------------------------------------------------
%  Scale of the challenge
% ------------------------------------------------------------------

\subsection{The Scale of the Challenge}
\label{sec:scale}

The fragmentation problem becomes combinatorially intractable when
viewed across the full landscape of computational imaging.
PWM's current registry includes 64 distinct modalities spanning five
physical carriers.
For each modality, at least 5 types of operator mismatch are
physically relevant (alignment, calibration drift, manufacturing
tolerance, environmental variation, model approximation error).
The solver space includes at least 10 competitive reconstruction
methods per modality.

The resulting evaluation space is therefore on the order of:
\begin{equation}
  64 \text{ modalities} \times 5 \text{ mismatch types}
  \times 10 \text{ solvers}
  = 3{,}200 \text{ experiments (minimum)}.
\end{equation}
With multiple severity levels per mismatch type and multiple test
scenes per experiment, the full evaluation matrix easily exceeds
$10^5$ individual runs.
No ad hoc, lab-by-lab evaluation effort can cover this space.
What is needed is not more careful experimentation but a fundamentally
different approach: \emph{standardized infrastructure} for evaluation.

% ------------------------------------------------------------------
%  Summary
% ------------------------------------------------------------------

\subsection{Summary: The Case for Infrastructure}
\label{sec:problem_summary}

The evidence presented in this section leads to three conclusions:

\begin{enumerate}
  \item \textbf{Neural solvers amplify mismatch.}
        The most capable architectures exhibit the largest sensitivity to
        operator error, with MST-L losing 16.72\,dB versus GAP-TV's
        4.56\,dB (\cref{tab:sensitivity}).

  \item \textbf{Evaluation fragmentation hides the problem.}
        Incomparable metrics, datasets, noise models, and mismatch
        definitions prevent the community from recognizing the
        severity of the sim-to-real gap.

  \item \textbf{The scale is beyond ad hoc effort.}
        With 64 modalities, 5+ mismatch types, and 10+ solvers, the
        evaluation space requires automated, standardized
        infrastructure---not heroic individual experiments.
\end{enumerate}

The field does not need a better solver.
It needs a better rail.
The next section presents the conceptual framework---drawn from
SolveEverything.org---that guides the design of that rail.
